% Direct-authority diplomatic French for printed p.164.
\SourcePage{164}

\LayerTarget{tohoku:diagram:3.4.2:H1-square}
\[
\begin{array}{ccc}
H^1(X,G)&\to&H^1(X,P)\\
\big\downarrow&&\big\downarrow\\[0.45em]
H^1(X,G)&\to&H^1(X,H)
\end{array}
\]
\LayerTarget{tohoku:eq:3.4.2:obstruction-cohomology}
où les flèches verticales sont des bijections) : et d'autre part
$H^1(X,{F'}^Q)$ est cannoniquement isomorphe à
$H^2(X,F^Q)=H^2(X,F^E)$, puisque de la suite exacte
\LayerTarget{tohoku:eq:3.4.2:twisted-flabby-sequence}
$0\to F\to\overline F\to F'\to0$
on déduit la suite exacte
$0\to F^Q\to\overline F^Q\to {F'}^Q\to0$
et que, $\overline F^Q$ étant flasque puisque localement isomorphe au faisceau
flasque $\overline F$, $H^i(X,\overline F^Q)=0$ pour $i>0$.
\LayerTarget{tohoku:eq:3.4.2:boundary-as-minus-dQ}
Il suffira donc de poser
$\partial(E)=-d(Q)\in H^2(X,F^E)$
pour satisfaire aux conditions voulues. Il reste donc à définir $d(Q)$, ayant
les propriétés de fonctoralité et ``d'exactitude'' voulues. C'est ce qui
va être fait sous les conditions plus générales qui suivent où les notations
sont légèrement changées).

\LayerTarget{tohoku:construction:crossed-homomorphism-general}
Soit $P$ un faisceau de groupes sur $X$, $A$ un faisceau de groupes sur lequel
$P$ opère (à gauche), l'opération définie par $p\in P$ étant désignée par
$\sigma(p)$. Soit $u$ un homomorphisme croisé de $P$ dans $A$, i. e.
un'homomorphisme de faisceaux tel que, sur chaque fibre $F(x)$, on ait
$u(e)=e$, $u(pp')=u(p)\sigma(p)\mathbin{\cdot}u(p')$. Alors le sous-faisceau
$G$ de $P$ image inverse de la section nulle de $A$ par $u$ est un
sous-faisceau de groupes, et deux éléments de $P$ ont même image dans $A$ si
et seulement si ils définissent la même classe à droite mod. $G$.
\LayerTarget{tohoku:eq:rho-action}
D'autre part, pour tout $x\in X$, $p\in P(x)$, $a\in A(x)$, posons
\[
\rho(p)a=u(p)(\sigma(p)a).
\]
Dire que $u$ est un homomosphisme croisé signifie précisément que la formule
précédente définit une représentation $\rho$ de $P$ par des germes
d'automorphismes du faisceau \textit{d'ensembles} $A$. D'ailleurs
l'homomorphisme de faisceaux d'ensembles
$A\times A\to A$ défini par le produit dans $A$, est compatible
avec les opérations de $P$ opérant respectivement par $\rho,\sigma,\rho$ :

\LayerTarget{tohoku:eq:rho-product-compatibility}
\[
\rho(p)(ab)=(\rho(p)a))\sigma(p)b).
\]
\LayerTarget{tohoku:construction:associated-rho-torsor}
On en conclut, pour tout fibré $E$ à faisceau structural $P$, un homomorphisme
des faisceaux d'ensembles associés
$A(\rho)^E\times A(\sigma)^E\to A(\rho)^E$,
et il est immédiat qu'ainsi $A(\rho)^E$ devient un faisceau d'ensembles sur
lequel le faisceau de groupes $A(\sigma)^E$ opère à droite, et plus précisément
$A(\rho)^E$ est un faisceau \textit{principal} (à droite) sous $A(\sigma)^E$.
\LayerTarget{tohoku:def:dE-obstruction}
Nous pouvons considérer sa classe
$c(A(\rho)^E)\in H^1(X,A(\sigma)^E)$,
que nous dénoterons aussi par $d(E)$. Son annulation est la condition nécessaire
et suffisante pour qu'il existe une section du faisceau $A(\rho)^E$. Notons
d'ailleurs que le monomorphisme $P/G\to A$ déduit de $u$ est
compatible avec les opérations de $P$, opérant sur $P/G$ de la façon canonique
et sur $A$ par $\rho$ ; d'où un monomorphisme naturel pour les faisceaux
associés :
\LayerTarget{tohoku:eq:associated-quotient-monomorphism}
$(P/G)^E\to A(\rho)^E$,
bijectif si et seulement si $u$ est surjectif.
\LayerTarget{tohoku:proof:3.4.2:closure}
Par suite, l'existence d'une section de $(P/G)^E$, qui est la condition
nécessaire et suffisante pour que la classe $c(E)\in H^1(X,P)$ de $E$ soit
dans l'image de $H^1(X,G)$, implique l'existence d'une section de $A(\rho)^E$
c'est-à-dire la nullité de $d(A)$, et la réciproque est vraie si $u$ est un
épimorphisme de $P$ sur $A$. Ces considérations achèvent donc la démonstration
de la prop. 3.4.2.

\LayerTarget{tohoku:remark:3.4.2:1:p1}
\indent\textsc{Remarques.} 1. On a mis le signe $-$ dans la formule
$\partial(E)=-d(Q)$,

\TohokuContinuation{165}
